\documentclass{article}
\usepackage[T1]{fontenc}
\usepackage{lmodern}
\usepackage{graphicx}
\usepackage{amsmath,amssymb}
\usepackage[a4paper,margin=2cm]{geometry}
\usepackage[absolute,overlay]{textpos}
\setlength{\TPHorizModule}{1mm}
\setlength{\TPVertModule}{1mm}
\usepackage[indonesian]{babel}
\usepackage{hyperref}
\setlength{\emergencystretch}{2em}
% unit: Halaman 39

\begin{document}

% === Halaman 39 (llm) ===

\begin{itemize}
    \item Dua pihak yang berkomunikasi menyepakati parameter data sebagai berikut:
    \begin{enumerate}
        \item Persamaan kurva eliptik \[ y^2 \equiv x^3 + ax + b \pmod{p} \]
        \begin{itemize}
            \item Nilai a dan b
            \item Bilangan prima p
        \end{itemize}
        \item Grup eliptik yang dihitung dari persamaan kurva eliptik
        \item Titik basis (\textit{base point}) B $(x_{\text{B}}, y_{\text{B}})$, dipilih dari grup eliptik untuk operasi kriptografi.
    \end{enumerate}
\end{itemize}

\begin{itemize}
    \item Setiap pengguna membangkitkan pasangan kunci publik dan kunci privat
    \begin{itemize}
        \item Kunci privat = integer x, dipilih dari selang $[1, p - 1]$
        \item Kunci publik = titik Q, adalah hasil kali antara x dan titik basis B: $\text{Q} = \text{x} \cdot \text{B}$
    \end{itemize}
\end{itemize}

\vfill
\raggedleft
\textcolor{red}{\small *) Sumber bahan: \textbf{Debdeep Mukhopadhyay}, \textit{Elliptic Curve Cryptography}, \\ Dept of Computer Sc and Engg IIT Madras}

\end{document}
